% from ega3-7-fr.tex
\oldpage[III]{187}
est nul par définition (M, IV, 1), donc $T_p$ est exact à gauche. Montrons
inversement que $b)$ implique $c)$. Si
\[
  Z_{p+1}(P_\bullet)=\operatorname{Ker}(P_{p+1}\longrightarrow P_p),
\]
on a une suite exacte
\[
  0\longrightarrow Z_{p+1}(P_\bullet)\longrightarrow P_{p+1}
  \longrightarrow Z'_p(P_\bullet)\longrightarrow0
\]
dans laquelle $P_{p+1}$ et $Z'_p(P_\bullet)$ sont plats, donc
$Z_{p+1}(P_\bullet)$ est plat ($0_{\mathrm I}$, 6.1.2). On va prendre
\[
  P'_i=P_i\quad\text{pour }i\ne p\text{ et }i\ne p+1,
  \qquad
  P'_p=Z'_p(P_\bullet),
  \qquad
  P'_{p+1}=Z_{p+1}(P_\bullet);
\]
pour différentielle $d'_i:P'_i\to P'_{i-1}$, on prendra celle du complexe
$P_\bullet$ pour $i\ne p$ et $i\ne p+1$, $0$ pour $i=p+1$ et pour $i=p$
l'homomorphisme $Z'_p(P_\bullet)\to P_{p-1}$ déduit de $d_p$ par passage au
quotient. Comme les $P_i$ sont plats, on a
\[
  Z'_i(P_\bullet\otimes M)=Z'_i(P_\bullet)\otimes M,
  \qquad
  Z_i(P_\bullet\otimes M)=Z_i(P_\bullet)\otimes M
\]
et
\[
  B_i(P_\bullet\otimes M)=B_i(P_\bullet)\otimes M
\]
(en posant $B_i(P_\bullet)=\operatorname{Im}(P_{i+1}\to P_i)$); on en conclut
aussitôt pour tout $M$ des isomorphismes fonctoriels
\[
  H_i(P_\bullet\otimes M)\xrightarrow{\sim}H_i(P'_\bullet\otimes M)
\]
pour tout $i$, et la vérification du fait qu'il s'agit d'un isomorphisme de
$\partial$-foncteurs découle sans peine de la définition de $\partial$
(M, IV, 1).
\end{proof}

On remarquera que les conditions de (7.4.2) impliquent aussi que
$B_p(P_\bullet)$ est plat, car on a une suite exacte
\[
  0\longrightarrow B_p(P_\bullet)\longrightarrow P_p
  \longrightarrow Z'_p(P_\bullet)\longrightarrow0,
\]
dans laquelle $P_p$ et $Z'_p(P_\bullet)$ sont plats ($0_{\mathrm I}$, 6.1.2).

\begin{corollary}[7.4.3]
Supposons que $A$ soit un anneau noethérien régulier de dimension $1$
(autrement dit un produit d'anneaux de Dedekind
($0_{\mathrm{IV}}$, 17.1.3 et 17.3.7), par exemple un anneau principal). Alors,
pour que $T_p$ soit exact à gauche, il faut et il suffit que $T_p(A)$ soit un
$A$-module plat. Pour que $T_p$ soit exact, il faut et il suffit que $T_p(A)$
et $T_{p-1}(A)$ soient des $A$-modules plats.
\end{corollary}

\begin{proof}
Rappelons que pour un module $M$ sur un anneau de Dedekind, il revient au même
de dire que $M$ est plat ou qu'il est sans torsion ($0_{\mathrm I}$, 6.3.3 et
6.3.4); sous les hypothèses de (7.4.3), tout sous-module d'un $A$-module plat
est donc plat.

La seconde assertion de (7.4.3) résulte de la première, puisque dire que $T_p$
est exact signifie que $T_p$ et $T_{p-1}$ sont exacts à gauche. Pour démontrer
la première assertion, notons que l'on a une suite exacte
\[
  0\longrightarrow H_p(P_\bullet)\longrightarrow Z'_p(P_\bullet)
  \longrightarrow B_{p-1}(P_\bullet)\longrightarrow0
\]
dans laquelle $B_{p-1}(P_\bullet)$ est un $A$-module plat, en tant que
sous-module du $A$-module plat $P_{p-1}$. Il est donc équivalent de dire que
$H_p(P_\bullet)$ est plat ou que $Z'_p(P_\bullet)$ est plat
($0_{\mathrm I}$, 6.1.2).
\end{proof}

Les applications les plus importantes de (7.4.2) sont les suivantes :

\begin{proposition}[7.4.4]
Soient $A$ un anneau noethérien, $P_\bullet$ un complexe de $A$-modules plats;
on suppose, soit que les $P_i$ sont de type fini, soit que les
$H_i(P_\bullet)$ sont des $A$-modules de type fini et qu'il existe $i_0$ tel
que $H_i(P_\bullet)=0$ pour $i<i_0$. Soit $T$ le foncteur homologique défini
par (7.4.1.1). Alors l'ensemble $U$ des
$y\in\operatorname{Spec}(A)$ tels que $(T_p)_y$ (7.1.4) soit exact à droite
(resp. exact à gauche, exact) est ouvert dans $\operatorname{Spec}(A)$.
\end{proposition}

\begin{proof}
Dans la seconde hypothèse sur $P_\bullet$, on peut remplacer $P_\bullet$ par
un complexe $P'_\bullet$ de
